\PhanI 
%\loai{Tìm độ biến dạng tại vị trí cân bằng }
\begin{ex}%[3P1B6-2][Loai1]%Bài 15 dạng 1
Một con lắc lò xo treo thẳng đứng gồm lò xo có độ cứng 50 N/m và vật nặng khối lượng 100 g đang dao động điều hòa. Lấy $g = 10\ m/s^2$. Tại vị trí cân bằng lò xo biến dạng một đoạn là
\choice
{5 cm}
{0,5 cm}
{\True 2 cm}
{2 mm}
\loigiai{
Ta có: $ \Delta \ell_0=\dfrac{mg}{k}= 2\ cm$.
}
%\haidong{4}
\end{ex} \haidong{20}
\begin{ex}%[3P1K6-2][Loai1] 
Một con lắc lò xo dao động điều hòa theo phương thẳng đứng, trong quá trình dao động của vật lò xo có chiều dài biến thiên từ 12 cm đến 20 cm. Biên độ dao động của vật là
\choice
{8 cm}
{\True 4 cm}
{16 cm}
{10 cm}
\loigiai{
Biên độ dao động của vật: $A=\dfrac{\ell_{\max}-\ell_{\min}}{2}=4\ cm$.
}
%\haidong{6}
\end{ex} \haidong{20}
%\loai{Tìm chu kì}

\begin{ex}%[3P1K6-2][Loai1]%Bài 15 dạng 1
Một con lắc lò xo treo thẳng đứng, dao động điều hòa dọc theo quỹ đạo dài 6 cm. Khi vật ở vị trí cao nhất, lò xo bị dãn 1 cm. Lấy $g = 10\ m/s^2,\ \pi^2 = 10$. Chu kì dao động của con lắc bằng
\choice
{0,5 s}
{0,6 s}
{\True 0,4 s}
{0,3 s}
\loigiai{
\begin{itemize}
\item Ở vị trí cao nhất (biên bên trên), lò xo bị dãn 1 cm
\item Vật ở dưới vị trí lò xo tự nhiên 1 cm $\Rightarrow \Delta \ell_0 = A + 1\  cm = 4\ cm \Rightarrow T = 0,4\ s$.
\end{itemize}
}
%\haidong{4}
\end{ex} \haidong{20}
%\loai{Tìm chiều dài max, min của lò xo}

%\loai{Tìm chiều dài lò xo ở vị trí cân bằng}
\begin{ex}%[3P1B6-2][Loai1]%Bài 15 dạng 1
Một con lắc lò xo treo thẳng đứng gồm lò xo có dài tự nhiên là 30 cm và vật nhỏ có khối lượng 200 g, lò xo có độ cứng 50 N/m. Lấy $g = 10\ m/s^2$. Chiều dài lò xo khi vật ở vị trí cân bằng là
\choice
{32 cm}
{\True 34 cm}
{35 cm}
{33 cm}
\loigiai{
Ta có: $ \Delta \ell_0=\dfrac{mg}{k}=  4\ cm \Rightarrow \ell cb = \ell_0 + \Delta \ell_0  = 34\ cm$.
}
%\haidong{4}
\end{ex} \haidong{20}

%\loai{Tìm chiều dài tự nhiên của lò xo}
\begin{ex}%[3P1K6-2][Loai1]%Bài 15 dạng 1
Con lắc lò xo treo thẳng đứng dao động với phương trình $x=8\sin \left( 20t+\dfrac{\pi}{2} \right)\ cm.$ Lấy $g = 10\ m/s^2$. Biết chiều dài lớn nhất của lò xo là 92,5 cm. Chiều dài tự nhiên của lò xo là
\choice
{\True 82 cm}
{84,5 cm}
{55 cm}
{61 cm}
\loigiai{
\begin{itemize}
\item Ta có: $ \Delta \ell_0=\dfrac{g}{\omega^2}=  2,5\ cm$.
\item Bài cho: $\ell_{\max} = \ell_0 + \Delta \ell_0  + A = 92,5\ cm \Rightarrow \ell_0 = 82\ cm$.
\end{itemize}
}
%\haidong{4}
\end{ex} \haidong{20}
\begin{ex}%[3P1B6-2][Loai1]%Bài 15 dạng 1
Một con lắc lò xo treo thẳng đứng dao động điều hòa với tần số 4,5 Hz. Trong quá trình dao động chiều dài lò xo biến thiên từ 40 cm đến 56 cm. Lấy $g = 10\ m/s^2$. Chiều dài tự nhiên của lò xo là
\choice
{48 cm}
{\True 46,75 cm}
{42 cm}
{40 cm}
\loigiai{
\begin{itemize}
\item Ta có: $f = 4,5\ Hz \Rightarrow \omega  = 9\pi \Rightarrow  \Delta \ell_0=\dfrac{g}{\omega^2}\approx 1,25\ cm$.
\item Lại có: $\ell_{cb}=\dfrac{\ell_{\max}+\ell_{\min}}{2}=  48\ cm = \ell_0+\Delta \ell_0 \Rightarrow \ell_0=  46,75\ cm$.
\end{itemize}
}
%\haidong{4}
\end{ex} \haidong{20}

%\loai{Biết chiều dài max, min tìm A, $\Delta \ell_0$ và năng lượng}


\begin{ex}%[3P1K6-2][Loai1] 
Một con lắc lò xo dao động điều hòa theo phương ngang. Trong quá trình dao động, chiều dài của lò xo biến thiên từ 22 cm đến 30 cm. Biết lò xo của con lắc có độ cứng 100 N/m. Khi vật cách vị trí biên 3 cm thì động năng của vật là
\choice
{0,045 J}
{0,0375 J}
{\True 0,075 J}
{0,035 J}
\loigiai{
\begin{itemize}
\item Biên độ dao động của con lắc là: $A=\dfrac{30-22}{2}=4\ cm$
\item Áp dụng định luật bảo toàn năng lượng ta có:
\begin{align*}
W_{\text{đ}}+W_t=\dfrac{kA^2}{2}\Rightarrow W_{\text{đ}}=\dfrac{kA^2}{2}-W_t=\dfrac{kA^2}{2}-\dfrac{kx^2}{2}=\dfrac{100.0,04^2}{2}-\dfrac{100.0,01^2}{2}=0,075\ J.
\end{align*}
\end{itemize}
}%<MyLT1>
%\haidong{3}
\end{ex} \haidong{20}
%\loai{Tìm độ dãn lớn nhất}
\begin{ex}%[3P1K6-2][Loai1] 
Con lắc lò xo treo thẳng đứng với đầu cố định của lò xo ở trên cao, quả nặng ở dưới thấp. Con lắc đang dao động điều hòa trên phương thẳng đứng, quanh vị trí cân bằng, với tần số bằng $\pi$ Hz. Biết độ nén lớn nhất của lò xo trong quá trình dao động là 2,5 cm. Cho gia tốc trọng trường $g = \pi^2 = 10\ m/s^2$. Độ giãn cực đại của lò xo trong lò xo trong quá trình dao động là
\choice
{10 cm}
{2,5 cm}
{\True 7,5 cm}
{2 cm}
\loigiai{
\begin{itemize}
\item Ta có: $\omega = 2\pi.f = 2\pi.\pi = 20\ rad/s$.
\item Tại vị trí cân bằng, lò xo dãn một đoạn $\Delta \ell_0 = \dfrac{g}{\omega^2} = \dfrac{10}{20^2} = 0,025\ m = 2,5\ cm$.
\item Độ nén lớn nhất của lò xo trong quá trình dao động là 2,5 cm $\Rightarrow A = +2,5 + 2,5 = 5\ cm$.
\item Độ giãn cực đại của lò xo trong lò xo trong quá trình dao động là:  $A + \Delta \ell_0=5 + 2,5=7,5\ cm$.
\end{itemize}
}
%\haidong{6}
\end{ex} \haidong{20}
%\loai{Tìm độ nén lớn nhất}

%\loai{Tìm chiều dài lò xo ở li độ và thời điểm cho trước}
\begin{ex}%[3P1K6-2][Loai1]%Bài 15 dạng 1
Con lắc lò xo treo thẳng đứng. Chọn chiều dương hướng thẳng đứng từ dưới lên trên. Khi vật dao động thì $\ell_{\max} = 100$ cm và $\ell_{\min} = 80\ cm$. Chiều dài của lò xo lúc vật ở li độ $x = -2$ cm là
\choice
{88 cm}
{82 cm}
{78 cm}
{\True 92 cm}
\loigiai{
\begin{itemize}
\item Ta có: $\ell_{cb}=\dfrac{\ell_{\max}+\ell_{\min}}{2}= 90\ cm$.
\item Chiều dương Ox hướng lên $\Rightarrow x = -2\ cm < 0$: vật đang ở dưới vị trí cân bằng đoạn 2 cm $\Rightarrow \ell  = \ell_{cb} + 2 = 92\ cm$.
\end{itemize}
}%<MyLT1>
%\haidong{8}
\end{ex} \haidong{20}
\begin{ex}%[3P1K6-2][Loai1]%Bài 15 dạng 1
Con lắc lò xo treo thẳng đứng, dao động điều hòa với phương trình $x=5\cos  \left( 4\pi t+\dfrac{\pi}{3} \right)\ cm$, chiều dương hướng xuống. Chiều dài tự nhiên của lò xo là 40 cm. Lấy $g =\pi^2= 10\ m/s^2$. Chiều dài lò xo tại thời điểm $\dfrac{1}{3}$ s là
\choice
{43,5 cm}
{\True 48,75 cm}
{43,75 cm}
{46,25 cm}
\loigiai{
\begin{itemize}
\item Ta có: $ \Delta \ell_0=\dfrac{g}{\omega^2}=  6,25\ cm;\ \Rightarrow \ell_{cb} = \ell_0 + \Delta \ell_0  = 46,25\ cm$.
\item Tại $t = \dfrac{1}{3}\ s \Rightarrow \phi \equiv -\dfrac{\pi}{3}: \  x=\dfrac{A}{2}=2,5\ cm$.
\item Chiều dương hướng xuống $\Rightarrow$ vật đang ở dưới vị trí cân bằng 2,5 cm $\Rightarrow \ell  = \ell_{cb} + 5\ cm = 48,75\ cm$.
\end{itemize}
}
%\haidong{7}
\end{ex} \haidong{20}
%\loai{Biết độ biến dạng và tốc độ, tìm các đại lượng}
\begin{ex}%[3P1K6-2][Loai1]%Bài 15 dạng 1
Một con lắc lò xo treo thẳng đứng gồm lò xo có độ cứng 100 N/m và vật nhỏ có khối lượng 100 g. Giữ vật theo phương thẳng đứng làm lò xo dãn 3 cm rồi truyền cho nó tốc độ $20\pi \sqrt{3}$ cm/s theo phương thẳng đứng, sau đó vật dao động điều hòa. Lấy $g = 10\ m/s^2,\ \pi^2=10$. Biên độ dao động của vật là
\choice
{5,46 cm}
{\True 4 cm}
{4,58 cm}
{2,54 cm}
\loigiai{
\begin{itemize}
\item Ta có: $\omega  = 10\pi\  rad/s$ và $\Delta \ell_0 = 1\ cm$.
\item Khi truyền tốc độ: lò xo đang dãn 3 cm $\Rightarrow$ vật ở dưới vị trí lò xo tự nhiên 3 cm $\Rightarrow$ vật đang cách vị trí cân bằng $\left|x\right| = \left|\Delta \ell_0 - 3\ cm\right| = 2\ cm$ và $\left|v\right| =20\pi \sqrt{3}\ cm/s \Rightarrow  A=\sqrt{x^2+\dfrac{v^2}{\omega^2}}=  4\ cm.$
\end{itemize}
}
%\haidong{9}
\end{ex} \haidong{20}
%\loai{Tìm tốc độ cực đại}

\begin{ex}%[3P1K6-2][Loai1] 
Con lắc lò xo treo thẳng đứng với đầu cố định của lò xo ở trên cao, quả nặng ở dưới thấp. Con lắc đang dao động điều hòa trên phương thẳng đứng, quanh vị trí cân bằng, với chu kì bằng $\dfrac{1}{\pi}$ s. Biết độ nén lớn nhất của lò xo trong quá trình dao động là 2,5 cm. Cho gia tốc trọng trường $g = \pi^2 = 10\ m/s^2$. Tốc độ cực đại của vật trong lò xo trong quá trình dao động là
\choice
{10 cm/s}
{25 cm/s}
{7,5 cm/s}
{\True 1 m/s}
\loigiai{
\begin{itemize}
\item Ta có: $\omega = \dfrac{2\pi}{T} = \dfrac{2\pi}{\dfrac{1}{\pi}} = 20\ rad/s$.
\item Tại vị trí cân bằng, lò xo dãn một đoạn $\Delta \ell_0 = \dfrac{g}{\omega^2} = \dfrac{10}{20^2} = 0,025\ m = 2,5\ cm$.
\item Mà độ nén lớn nhất của lò xo trong quá trình dao động là 2,5 cm $\Rightarrow A = 2,5 + 2,5 = 5\ cm\Rightarrow v_{\max} = \omega.A = 20.5 = 100\ cm/s = 1\ m/s$.
\end{itemize}
}
%\haidong{12}
\end{ex} \haidong{20}

%\loai{Tìm tốc độ tại vị trí bất kì}

%\loai{Viết phương trình dao động}


























%\loai{Tính lực đàn hồi cực đại,cực tiểu}

\begin{ex}%[3P1K7-2][Loai1]
Một con lắc lò xo treo thẳng đứng gồm vật nặng có khối lượng 1 kg và lò xo khối lượng không đáng kể. Giữ vật ở vị trí dưới vị trí cân bằng sao cho khi đó lực đàn hồi tác dụng lên vật có độ lớn 12 N rồi thả nhẹ cho vật dao động điều hòa. Lấy $g = 10\ m/s^2$. Lực đàn hồi của lò xo có độ lớn nhỏ nhất trong quá trình vật dao động là
\choice
{0}
{4 N}
{\True 8 N}
{22 N}
\loigiai{
\begin{itemize}
\item Đoạn kéo xuống chính là biên độ (gọi là A).
\item Khi đó độ lớn lực đàn hồi $\left|F\right| = 12N = k(\Delta \ell + A) = 10N + kA$ $\Rightarrow$ kA = 2N
\item Độ lớn lực đàn hồi nhỏ nhất là $\left|F\right|_{\min} = k(\Delta \ell - A) = 8$ N (xảy ra ở biên trên).
\end{itemize}
}
%\haidong{5}
\end{ex} \haidong{20}


%\loai{Tính tỉ số lực đàn hồi}

\begin{ex}%[3P1K7-2][Loai1]
Con lắc lò xo treo thẳng đứng dao động với phương trình $x=12\cos  \left( 10t+\dfrac{\pi}{3} \right)cm$ tại nơi có $g = 10$ $m/s^2$. Tỉ số độ lớn lực đàn hồi khi vật ở biên dưới và biên trên là
\choice
{3}
{8}
{\True 11}
{12}
\loigiai{
$\omega  = 10$ rad/s $\Rightarrow$ $\Delta \ell = 0,1$ m $\Rightarrow$ $\dfrac{{{\left| F \right|}_{\max}}}{{{\left| F \right|}_{\min}}}=\dfrac{\Delta \ell+A}{\Delta \ell-A}=11$}

%\haidong{5}
\end{ex} \haidong{20}

%\loai{Tính lực đàn hồi tại vị trí bất kì}
\begin{ex}%[3P1K7-2][Loai1]
Con lắc lò xo treo thẳng đứng đang dao động có độ dãn lò xo khi vật ở vị trí cân bằng là 10 cm. Vật nặng dao động trên chiều dài quỹ đạo là 24 cm. Lò xo có độ cứng 40 N/m. Độ lớn lực tác dụng vào điểm treo khi lò xo có chiều dài ngắn nhất là
\choice
{\True 0,8 N}
{8 N}
{80 N}
{5,6 N}
\loigiai{
\begin{itemize}
\item $A = 0,12 m > 0,1 m = \Delta \ell$
\item Khi lò xo ngắn nhất (biên trên), lò xo bị nén: ($A - \Delta \ell$) = 0,02m
\item $\Rightarrow$ Lò xo tác dụng lực đẩy lên điểm treo với độ lớn: $\left|F\right| = k(A - \Delta \ell) = 0,8$ N.
\end{itemize}
}
%\haidong{5}
\end{ex} \haidong{20}

%\loai{Từ tỉ số lực đàn hồi tìm các đại lượng}


%\loai{Từ lực đàn hồi max (hoặc min), tìm các đại lượng}
\begin{ex}%[3P1G7-2][Loai1] 
Con lắc lò xo dao động điều hòa theo phương thẳng đứng có năng lượng dao động $W=2.10^{-2}$ J và lực đàn hồi cực đại của lò xo $F_{max}=4\ N$ . Lực đàn hồi của lò xo khi vật ở vị trí cân bằng là $F = 2$ N. Biên độ dao động của vật là
\choice
{\True 2 cm}
{4 cm}
{5 cm}
{3 cm}
\loigiai{
Theo đề bài ta có:
$F = k.\Delta \ell = 2$; 
$F_{max} = k\left(\Delta \ell + A\right) = k\Delta \ell + kA = 2 + kA = 4 \Rightarrow kA = 2$
$\Rightarrow W = \dfrac{1}{2}k.A^2 \Rightarrow A=2\ cm$.
}
%\haidong{6}
\end{ex} \haidong{20}

%\loai{Từ lực đàn hồi max và min, tìm các đại lượng}



%\loai{Từ lực đàn hồi tại vị trí bất kì tìm các đại lượng}

%



















%\loai{Tìm thời gian, thời điểm theo độ biến dạng}
\begin{ex}%[3P1KA-1][Loai1] 
Một con lắc lò xo có khối lượng $m=100$ g và lò xo có độ cứng $k=100$ N/m, dao động trên mặt phẳng nằm ngang. Kéo vật khỏi vị trí cân bằng một khoảng 3 cm rồi truyền cho vật vận tốc bằng $30\pi \sqrt{3}$ cm/s theo chiều hướng ra xa vị trí cân bằng để vật bắt đầu dao động điều hoà, chọn gốc thời gian lúc vật bắt đầu dao động, lấy $\pi ^2=10$. Khoảng thời gian ngắn nhất kể từ khi vật bắt đầu dao động điều hoà đến khi lò xo bị nén cực đại là
\choice
{$\dfrac{3}{20}$ s}
{$\dfrac{3}{10}$ s}
{\True $\dfrac{2}{15}$ s}
{$\dfrac{1}{15}$ s}
\loigiai{
\begin{itemize}
\item Ta có $\omega = \sqrt{\dfrac{k}{m}} = 10\pi\ rad/s$.
\item Chọn hệ trục x'Ox gốc O trùng vị trí cân bằng chiều dương hướng ra xa vị trí cân bằng.
\item Biên độ $A = \sqrt{x^2 + \dfrac{v^2}{\omega^2}} = 6\ cm$.
\item Ban đầu vật ở li độ $x=3\ cm$ theo chiều dương.
\item Lò xo nằm ngang nên lò xo bị nén cực đại ở vị trí $x=-A$.
\item Như vậy ta có thời gian từ lúc bắt đầu dao động tới lúc lò xo bị nén cực đại là: $t = \dfrac{T}{6} + \dfrac{T}{2} = \dfrac{2T}{3} = \dfrac{4\pi}{3\omega} = \dfrac{2}{15}\ s$.
\end{itemize}
}
%\haidong{10}
\end{ex} \haidong{20}
%\loai{Tìm khoảng thời gian giữa hai lực đàn hồi cho trước}


%\loai{Tìm thời gian thoả mãn một lực đàn hồi cho trước}



%\loai{Loại khác}









%\loai{Khoảng thời gian lò xo bị nén trong một chu kì}
\begin{ex}%[3P1KA-2][Loai1] 
Một con lắc gồm lò xo được treo trên phương thẳng đứng với đầu cố định của lò xo ở trên cao, vật nhỏ ở dưới thấp. Lấy gần đúng gia tốc trọng trường $g=\pi ^2=10\ m/s^2$. Kích thích để vật nhỏ dao động theo phương thẳng đứng với tần số 2,5 Hz và tốc độ cực đại trong quá trình dao động là $40\pi$ cm/s. Trong một chu kì dao động, quãng thời gian lò xo bị nén là
\choice
{$\dfrac{1}{6}$ s}
{$\dfrac{1}{4}$ s}
{$\dfrac{1}{3}$ s}
{\True $\dfrac{2}{15}$ s}
\loigiai{
\immini{\begin{itemize}
\item Ta có: $\omega = 2\pi f = 5\pi\ rad/s \Rightarrow A = \dfrac{40\pi}{5\pi} = 8\ cm$.
\item Lại có: $\Delta \ell_0 = \dfrac{g}{\omega^2} = 0,04\ m = 4\ cm$.
\item Trong 1 chu kì, lo xo bị nén khi đi từ vị trí $M_1$ qua biên âm, tới vị trí $M_2$ như hình vẽ.\\
$\Rightarrow \Delta \varphi = \dfrac{\pi}{3} + \dfrac{\pi}{3} = \dfrac{2\pi}{3}$.\\
$t = \dfrac{T}{3} = \dfrac{1}{3}.\dfrac{2\pi}{\omega} = \dfrac{1}{3}.\dfrac{2\pi}{5\pi} = \dfrac{2}{15}\ s$.
\end{itemize}}{\begin{tikzpicture}[scale=1,thick,>=stealth']
\tkzInit[ymin=-3,ymax=3,xmin=-3,xmax=3]\tkzClip
\tkzDefPoints{-2.5/0/m, 2.5/0/n, 0/0/o, 0/2.5/p, 0/-2.5/q}
\tkzDefShiftPoint[o](120:2){1}
\tkzDefShiftPoint[o](-120:2){2}
\tkzDefPointBy[projection=onto m--n](1) \tkzGetPoint{h}
\draw(o)circle(2cm);
\tkzDrawSegments[dashed](1,2)
\tkzDrawSegments(p,q)
\tkzDrawSegments[->](m,n)
\tkzDrawSegments[blue,->](o,1 o,2)
\tkzDrawPoints[size=3](o,1,h,2)
\tkzMarkAngles[size=0.7cm,arrows=->](1,o,2)
\node at (1)[above left]{$M_1$};
\node at (2)[below left]{$M_2$};
\node at (2,0)[below right]{$8$};
\node at (h)[below left]{$-4$};
\end{tikzpicture}}
}
%\haidong{7}
\end{ex} \haidong{20}
%\loai{Tính tỉ số thời gian lò xo dãn - nén}


%\loai{Tính thời gian lực đàn hồi ngược chiều với lực kéo về}























\PhanII 
\setcounter{ex}{0}
\begin{ex}
Một con lắc lò xo có độ cứng $k$, vật nặng có khối lượng $m$ được kích thích dao động điều hòa theo phương thẳng đứng.
\choiceTF[t]
{Vị trí cân bằng của dao động cũng là vị trí lò xo có chiều dài tự nhiên}
{\True Tần số góc dao động là $\omega=\sqrt{\dfrac{k}{m}}$}
{\True Vị trí cân bằng của vật cách vị trí lò xo có chiều dài tự nhiên một khoảng $\Delta \ell_0=\dfrac{mg}{k}$}
{\True Khoảng thời gian ngắn nhất vật đi từ vị trí thấp nhất đến vị trí cao nhất là $t=\dfrac{T}{2}$}
\loigiai{
\begin{itemchoice}
\itemch Con lắc lò xo dao động thẳng đứng, vị trí cân bằng của dao động không trùng với vị trí lò xo có chiều dài tự nhiên.
\itemch
\itemch
\itemch
\end{itemchoice}
}
%\haidong{2}
\end{ex} \haidong{20}
\begin{ex}
Một vật có khối lượng $m=100 \ g$ được treo vào lò xo đặt thẳng đứng có độ cứng $k=40 \ N/m$. Khi vật đang ở vị trí cân bằng người ta truyền cho vật một vận tốc $v=80 \ cm/s$ hướng xuống dưới. Lấy $g=10 \ m/s^2$.
\choiceTF[t]
{\True Tần số góc dao động là $\omega=20\ rad/s$}
{\True Chiều dài quỹ đạo là $8 \ cm$}
{Độ lớn lực đàn hồi cực đại là $1,6 \ N$}
{Quy ước chiều dương hướng xuống, phương trình dao động là $x=8 \cos \left(20 t-\dfrac{\pi}{2}\right)\ cm$}
\loigiai{
\begin{itemchoice}
\itemch Tần số góc dao động là $\omega=\sqrt{\dfrac{k}{m}}=20\ rad/s$.
\itemch Chiều dài quỹ đạo là $L=2 \ A=2 \dfrac{v_{\max}}{\omega}=8 \ cm$
\itemch Độ lớn lực đàn hồi cực đại là $F_{dh}=k\left(\Delta \ell_0+A\right)=k\left(\dfrac{g}{\omega^2}+A\right)=2,6 \ N$.
\itemch Phương trình dao động là $x=4 \cos \left(20 t-\dfrac{\pi}{2}\right)\ cm$.
\end{itemchoice}
}
%\haidong{5}
\end{ex} \haidong{20}


















\begin{ex}
Một con lắc lò xo có độ cứng $k$ được gắn thêm vật có khối lượng $m$ được treo thẳng đứng. Biết khi vật ở vị trí cân bằng thì lò xo giãn một đoạn $\Delta \ell_0 =4\ cm$. Kích thích cho vật dao động điều hòa. Biết quỹ đạo dao động có chiều dài $16 \ cm$. Lấy $g=\pi^2=10 \ m/s^2$.
\choiceTF[t]
{Tần số góc có thể được tính theo công thức $\omega=\sqrt{\dfrac{\Delta \ell_0}{g}}$}
{\True Chu kì dao động là $T=0,4\ s$}
{Tốc độ cực đại trong quá trình dao động là $v=1,6\ m/s$}
{\True Trong một chu kì, thời gian lò xo giãn gấp đôi thời gian lò xo nén}
\loigiai{
\begin{itemchoice}
\itemch Tần số góc có thể được tính theo công thức $\omega=\sqrt{\dfrac{g}{\Delta \ell_0}}=5 \sqrt{10} rad/s$.
\itemch Chu kì dao động là $T=2 \pi \sqrt{\dfrac{\Delta \ell_0}{\ g}}=0,4\ s$.
\itemch Tốc độ cực đại trong quá trình dao động là $v=\omega A=1,26(\ m/s)$.
\itemch Thời gian lò xo nén là $t_{\text {nén}}=\dfrac{\alpha_{\text {nen}}}{\omega}$ với $\alpha_{\text {nen}}=2 \operatorname{arc} \cos \left(\dfrac{-\Delta \ell_0}{\ A}\right)=\dfrac{2 \pi}{3}$ và $\alpha_{\text {dan}}=2 \pi-\dfrac{2 \pi}{3}=\dfrac{4 \pi}{3}=2 \alpha_{\text {nen}}$ nên $t_{\text {dan}}=2 t_{\text {nen}}$
\end{itemchoice}
}
%\haidong{4}
\end{ex} \haidong{20}
\begin{ex}
Một vật có khối lượng $m=100 \ g$ được gắn vào lò xo có độ cứng $k$ được treo theo phương thẳng đứng. Kích thích cho vật dao động điều hòa. Biết chu kì dao động là $T=1\ s$. Lấy $g=\pi^2=10 \ m/s^2$.
\choiceTF[t]
{Khi vật đi đến vị trí lò xo có chiều dài tự nhiên thì vận tốc vật lớn nhất}
{Độ cứng lò xo $k=40\ N/m$}
{\True Khi vật ở vị trí cân bằng thì lò xo giãn một đoạn $25\ cm$}
{\True Nếu biên độ dao động của vật là $10 \ cm$ thì lò xo luôn giãn trong quá trình dao động}
\loigiai{
\begin{itemchoice}
\itemch Vận có vận tốc lớn nhất khi đi qua vị trí cân bằng.
\itemch Độ cứng lò xo $T=2 \pi \sqrt{\dfrac{m}{k}} \Rightarrow k=3,95(\ N/m)$.
\itemch Khi vật ở vị trí cân bằng thì lò xo giãn một đoạn $\Delta \ell_0=\dfrac{mg}{k}=25\ cm$.
\itemch Nếu biên độ dao động của vật là $10 \ cm$ thì $A<\Delta \ell_0$ nên lò xo luôn giãn trong quá trình dao động.
\end{itemchoice}
}
%\haidong{5}
\end{ex} \haidong{20}








\PhanIII 
\setcounter{ex}{0}

\begin{ex}
Con lắc lò xo treo thẳng đứng, dao động điều hòa với phương trình $x=5 \cos \left(4 \pi t+\dfrac{\pi}{3}\right)\ cm$. Chiều dài tự nhiên của lò xo là $40 \ cm$. Chọn chiều dương hướng lên. Lấy $g=\pi^2 \ m/s^2$. Chiều dài của lò xo là bao nhiêu centimet khi vật dao động được $\dfrac{2}{3} \ T$, kể từ thời điểm $t=0$ (làm tròn kết quả đến chữ số hàng phần mười)?
\shortans[oly]{43,8}
\loigiai{
43,8
}
%\haidong{5}
\end{ex} \haidong{20}

\begin{ex}
Một con lắc lò xo dao động theo phương thẳng đứng, $k=40 \ N/m;\ m=100 \ g$. Giữ vật theo phương thẳng đứng làm lò xo giãn $3,5 \ cm$ rồi truyền cho nó vận tốc $20 \ cm/s$ hướng lên để vật dao động điều hòa, lấy $g=10 \ m/s^2$. Biên độ dao động của vật là bao nhiêu centimet (làm tròn kết quả đến chữ số hàng phần trăm)?
\shortans[oly]{1,41}
\loigiai{
1,41
}
%\haidong{4}
\end{ex} \haidong{20}
\begin{ex}
Lò xo có chiều dài tự nhiên $\ell_0=60 \ cm$ treo thẳng đứng dao động với phương trình $x=4 \cos \left(10 t+\dfrac{\pi}{3}\right)\ cm$. Chọn chiều dương hướng lên và lấy $g=10 \ m/s^2$. Chiều dài lò xo ở thời điểm $t=0,75 \ T$ là bao nhiêu centimet (làm tròn kết quả đến chữ số hàng phần mười)?
\shortans[oly]{63,5}
\loigiai{
63,5
}
%\haidong{4}
\end{ex} \haidong{20}
\begin{ex}
Con lắc lò xo treo thẳng đứng dao động với phương trình $x=8 \sin \left(20 t+\dfrac{\pi}{2}\right)\ cm$. Lấy $g=10 \ m/s^2$. Biết chiều dài lớn nhất của lò xo là $92,5 \ cm$. Chiều dài tự nhiên của lò xo là bao nhiêu centimet?
\shortans[oly]{82}
\loigiai{
82
}
%\haidong{4}
\end{ex} \haidong{20}
\begin{ex}
Một con lắc lò xo dao động theo phương thẳng đứng, $k=62,5 \ N/m; m=100 \ g$. Giữ vật theo phương thẳng đứng làm lò xo giãn $3,2 \ cm$ rồi truyền cho nó vận tốc $60 \ cm/s$ hướng lên để vật dao động điều hòa. Biên độ dao động của vật là bao nhiêu centimet (làm tròn kết quả đến chữ số hàng phần trăm)?
\shortans[oly]{2,88}
\loigiai{
2,88
}
%\haidong{4}
\end{ex} \haidong{20}
\begin{ex} %[3P1KA-1][Loai1] 
Một con lắc lò xo dao động theo phương ngang. Lực đàn hồi cực đại tác dụng vào vật là 12 N. Khoảng thời gian ngắn nhất giữa hai lần liên tiếp vật chịu tác dụng của lực kéo lò xo $6\sqrt{3}\ N$ là 0,1 s. Chu kì dao động của vật là bao nhiêu giây?
\shortans[oly]{0,6}
\loigiai{
\begin{itemize}
\item Ta có: $\begin{cases}
F_1=kx_1 \\ 
F_{\max}=kA \\ 
\end{cases} \Rightarrow \dfrac{6\sqrt{3}}{12}=\dfrac{F_1}{F_{\max}}=\dfrac{x_1}{A}\Rightarrow x_1=\dfrac{A\sqrt{3}}{2}$. 
\item Vì lực kéo nên lúc ấy lò xo bị dãn $\Rightarrow $ Vật đi xung quanh vị trí biên từ $x=\dfrac{A\sqrt{3}}{2}$ đến $x=A$ rồi đến $x=\dfrac{A\sqrt{3}}{2}$.
\item Thời gian đi sẽ là $\Delta t=\dfrac{T}{12}+\dfrac{T}{12}=\dfrac{T}{6}=0,1\Rightarrow T=0,6\ s$.
\end{itemize}
}
%\haidong{6}
\end{ex} \haidong{20}